\chapter*{ضمیمهٔ ۸: تقابل‌های موضوعی}
\addcontentsline{toc}{chapter}{ضمیمهٔ ۸: تقابل‌های موضوعی}
\markboth{ضمیمهٔ ۸: تقابل‌های موضوعی}{ضمیمهٔ ۸: تقابل‌های موضوعی}

\begin{center}
\textit{تصویرسازی بصری از تنش‌های بنیادین اندیشهٔ سیاسی}
\end{center}

\section*{۱. فرد در برابر جمع}

\begin{center}
\resizebox{\columnwidth}{!}{%
\begin{tikzpicture}[scale=1, every node/.style={font=\small}]

% Spectrum line
\draw[line width=4pt, left color=rightblue, right color=leftred] (0,0) -- (14,0);

% Position markers
\foreach \x/\label in {
    1/{\rl{آنارکو-کاپیتالیسم}},
    3.5/{\rl{لیبرتارینیسم}},
    6/{\rl{لیبرالیسم}},
    8.5/{\rl{سوسیال دموکراسی}},
    11/{\rl{سوسیالیسم}},
    13/{\rl{کمونیسم}}
} {
    \fill[color=darkgray] (\x,0) circle (4pt);
    \node[below, align=center] at (\x,-0.5) {\label};
}

% Axis labels
\node[above, font=\bfseries, color=rightblue] at (1,0.8) {\rl{فرد مطلق}};
\node[above, font=\bfseries, color=leftred] at (13,0.8) {\rl{جمع مطلق}};

% Description boxes
\node[draw=rightblue, fill=rightblue!10, rounded corners, align=right, text width=5cm] at (2,3) {
\textbf{\rl{تأکید بر فرد:}}\\
\rl{• حقوق فردی مقدس}\\
\rl{• مالکیت خصوصی}\\
\rl{• مسئولیت شخصی}\\
\rl{• رقابت}
};

\node[draw=leftred, fill=leftred!10, rounded corners, align=right, text width=5cm] at (12,3) {
\textbf{\rl{تأکید بر جمع:}}\\
\rl{• منافع اجتماعی مقدم}\\
\rl{• مالکیت اجتماعی}\\
\rl{• مسئولیت جمعی}\\
\rl{• همکاری}
};

\end{tikzpicture}%
}
\end{center}

\section*{۲. آزادی در برابر برابری}

\begin{center}
\resizebox{\columnwidth}{!}{%
\begin{tikzpicture}[scale=0.9, every node/.style={font=\small}]

% Two-dimensional chart
\draw[line width=2pt, ->] (0,0) -- (10,0) node[right] {\rl{برابری}};
\draw[line width=2pt, ->] (0,0) -- (0,8) node[above] {\rl{آزادی}};

% Quadrant labels
\node[align=center] at (2.5,6) {\rl{لیبرتارینیسم}\\{\footnotesize\rl{آزادی بالا، برابری پایین}}};
\node[align=center] at (7.5,6) {\rl{لیبرالیسم اجتماعی}\\{\footnotesize\rl{آزادی بالا، برابری بالا}}};
\node[align=center] at (2.5,2) {\rl{اقتدارگرایی راست}\\{\footnotesize\rl{آزادی پایین، برابری پایین}}};
\node[align=center] at (7.5,2) {\rl{اقتدارگرایی چپ}\\{\footnotesize\rl{آزادی پایین، برابری بالا}}};

% Position dots for ideologies
\fill[color=rightblue] (1.5,7) circle (5pt) node[right] {\footnotesize\rl{ لیبرتارین}};
\fill[color=centergreen] (5,5.5) circle (5pt) node[right] {\footnotesize\rl{ لیبرال}};
\fill[color=leftred] (8,4) circle (5pt) node[right] {\footnotesize\rl{ سوسیال دموکرات}};
\fill[color=leftred!70] (7,1.5) circle (5pt) node[right] {\footnotesize\rl{ استالینیسم}};
\fill[color=darkgray] (2,1.5) circle (5pt) node[right] {\footnotesize\rl{ فاشیسم}};

% Ideal point?
\fill[color=goldaccent] (6,6) circle (6pt);
\node[above right, color=goldaccent] at (6,6) {\rl{آرمان؟}};

\end{tikzpicture}%
}
\end{center}

\section*{۳. مالکیت خصوصی در برابر مالکیت عمومی}

\begin{center}
\resizebox{\columnwidth}{!}{%
\begin{tikzpicture}[scale=0.85, every node/.style={font=\small}]

% Visual representation
\node[font=\bfseries\large] at (7,7) {\rl{طیف مالکیت}};

% Left side - Private
\node[draw=rightblue, fill=rightblue!20, rounded corners, minimum width=3cm, minimum height=1.5cm, align=center] at (1,4) {
\textbf{\rl{مالکیت خصوصی کامل}}\\
\rl{آنارکو-کاپیتالیسم}
};

% Middle-left
\node[draw=rightblue!70, fill=rightblue!10, rounded corners, minimum width=3cm, minimum height=1.5cm, align=center] at (4.5,4) {
\textbf{\rl{خصوصی با تنظیم}}\\
\rl{لیبرالیسم}
};

% Middle
\node[draw=centergreen, fill=centergreen!20, rounded corners, minimum width=3cm, minimum height=1.5cm, align=center] at (8,4) {
\textbf{\rl{اقتصاد مختلط}}\\
\rl{سوسیال دموکراسی}
};

% Middle-right
\node[draw=leftred!70, fill=leftred!10, rounded corners, minimum width=3cm, minimum height=1.5cm, align=center] at (11.5,4) {
\textbf{\rl{عمومی غالب}}\\
\rl{سوسیالیسم}
};

% Right side - Public
\node[draw=leftred, fill=leftred!20, rounded corners, minimum width=3cm, minimum height=1.5cm, align=center] at (15,4) {
\textbf{\rl{مالکیت عمومی کامل}}\\
\rl{کمونیسم}
};

% Examples below
\node[below, align=center] at (1,2.5) {\footnotesize\rl{همه چیز خصوصی}\\{\scriptsize\rl{جاده، پلیس، دادگاه}}};
\node[below, align=center] at (8,2.5) {\footnotesize\rl{خصوصی + دولتی}\\{\scriptsize\rl{بازار + دولت رفاه}}};
\node[below, align=center] at (15,2.5) {\footnotesize\rl{همه چیز عمومی}\\{\scriptsize\rl{کارخانه، زمین، خدمات}}};

% Arrow
\draw[<->, line width=2pt, color=darkgray] (1,1.5) -- (15,1.5);
\node[below] at (8,1.2) {\rl{میزان مالکیت دولتی/اجتماعی}};

\end{tikzpicture}%
}
\end{center}

\section*{۴. رقابت در برابر همکاری}

\begin{center}
\resizebox{\columnwidth}{!}{%
\begin{tikzpicture}[scale=0.9, every node/.style={font=\small}]

% Two columns comparison
\node[draw=rightblue, fill=rightblue!15, rounded corners, minimum width=6cm, minimum height=6cm, align=center] at (3,3) {};
\node[font=\bfseries, color=rightblue] at (3,5.5) {\large\rl{رقابت}};

\node[draw=leftred, fill=leftred!15, rounded corners, minimum width=6cm, minimum height=6cm, align=center] at (11,3) {};
\node[font=\bfseries, color=leftred] at (11,5.5) {\large\rl{همکاری}};

% Competition side
\node[align=right, text width=5cm] at (3,3.5) {
\rl{• انگیزه برای نوآوری}\\
\rl{• کارایی اقتصادی}\\
\rl{• انتخاب مصرف‌کننده}\\
\rl{• پاداش استعداد}
};

\node[align=right, text width=5cm, color=leftred!70] at (3,1) {
\rl{⚠ نابرابری}\\
\rl{⚠ انحصار}\\
\rl{⚠ استثمار}
};

% Cooperation side
\node[align=right, text width=5cm] at (11,3.5) {
\rl{• همبستگی اجتماعی}\\
\rl{• کاهش نابرابری}\\
\rl{• تأمین همگانی}\\
\rl{• هدف مشترک}
};

\node[align=right, text width=5cm, color=rightblue!70] at (11,1) {
\rl{⚠ کاهش انگیزه؟}\\
\rl{⚠ ناکارآمدی؟}\\
\rl{⚠ بوروکراسی}
};

% Middle - synthesis?
\draw[<->, line width=3pt, color=centergreen] (6,3) -- (8,3);
\node[above, color=centergreen, font=\bfseries] at (7,3.3) {\rl{ترکیب؟}};

\end{tikzpicture}%
}
\end{center}

\section*{۵. آزادی منفی در برابر آزادی مثبت}

\begin{tcolorbox}[
    enhanced,
    colback=lightbg,
    colframe=darkgray,
    boxrule=2pt,
    arc=8pt,
    sidebyside,
    sidebyside align=top,
    lefthand width=0.45\textwidth
]
\begin{center}
\textbf{\Large\rl{آزادی منفی}}\\[3mm]
\textbf{\rl{«آزادی از...»}}
\end{center}

\begin{itemize}
    \item آزادی از دخالت دیگران
    \item آزادی از اجبار دولت
    \item «رهایم بگذارید»
    \item حقوق منفی: حق نداشتن مزاحم
\end{itemize}

\textbf{\rl{نمایندگان:}} آیزایا برلین، هایک، لیبرتارین‌ها

\textbf{\rl{نقد:}} آزادی صوری؛ گرسنه آزاد است؟

\tcblower

\begin{center}
\textbf{\Large\rl{آزادی مثبت}}\\[3mm]
\textbf{\rl{«آزادی برای...»}}
\end{center}

\begin{itemize}
    \item آزادی برای تحقق خود
    \item توانمندسازی واقعی
    \item «کمکم کنید بتوانم»
    \item حقوق مثبت: حق داشتن منابع
\end{itemize}

\textbf{\rl{نمایندگان:}} روسو، مارکس، سوسیال دموکرات‌ها

\textbf{\rl{نقد:}} بهانه برای پدرسالاری دولتی؟
\end{tcolorbox}

\section*{۶. کوتاه‌مدت در برابر بلندمدت}

\begin{center}
\resizebox{\columnwidth}{!}{%
\begin{tikzpicture}[scale=0.9, every node/.style={font=\small}]

% Timeline arrow
\draw[line width=3pt, ->] (0,3) -- (14,3);
\node[above] at (2,3.2) {\rl{امروز}};
\node[above] at (12,3.2) {\rl{آینده}};

% Short-term focused
\node[draw=goldaccent, fill=goldaccent!20, rounded corners, align=center] at (3,5) {
\textbf{\rl{تمرکز کوتاه‌مدت}}\\[2mm]
\rl{• نتایج فوری}\\
\rl{• انتخابات بعدی}\\
\rl{• سود سه‌ماهه}
};

% Long-term focused
\node[draw=centergreen, fill=centergreen!20, rounded corners, align=center] at (11,5) {
\textbf{\rl{تمرکز بلندمدت}}\\[2mm]
\rl{• نسل‌های آینده}\\
\rl{• پایداری}\\
\rl{• سرمایه‌گذاری}
};

% Issues
\node[below, align=center] at (3,1.5) {
\rl{بحران اقلیم: «بعداً»}\\
\rl{بدهی ملی: «مشکل بعدی‌ها»}
};

\node[below, align=center] at (11,1.5) {
\rl{بحران اقلیم: «الان»}\\
\rl{بدهی: «مسئولیت ما»}
};

% Tension
\draw[<->, line width=2pt, dashed, color=leftred] (5,5) -- (9,5);
\node[above, color=leftred] at (7,5.2) {\rl{تنش}};

\end{tikzpicture}%
}
\end{center}

\section*{۷. نظم در برابر آزادی}

\begin{center}
\resizebox{\columnwidth}{!}{%
\begin{tikzpicture}[scale=0.9, every node/.style={font=\small}]

% Main spectrum
\draw[line width=4pt, left color=darkgray, right color=centergreen] (0,0) -- (14,0);

% Labels
\node[above, font=\bfseries] at (1,0.5) {\rl{نظم مطلق}};
\node[above, font=\bfseries] at (13,0.5) {\rl{آزادی مطلق}};

% Position markers
\fill[color=darkgray] (1,0) circle (5pt);
\node[below] at (1,-0.3) {\rl{توتالیتاریسم}};

\fill[color=rightblue] (4,0) circle (5pt);
\node[below] at (4,-0.3) {\rl{اقتدارگرایی}};

\fill[color=centergreen] (7,0) circle (5pt);
\node[below] at (7,-0.3) {\rl{دموکراسی لیبرال}};

\fill[color=goldaccent] (10,0) circle (5pt);
\node[below] at (10,-0.3) {\rl{لیبرتارینیسم}};

\fill[color=violet] (13,0) circle (5pt);
\node[below] at (13,-0.3) {\rl{آنارشیسم}};

% Optimal zone?
\draw[line width=2pt, color=centergreen, decorate, decoration={brace, amplitude=10pt, mirror}] (5,-1) -- (9,-1);
\node[below, color=centergreen] at (7,-1.5) {\rl{محدودهٔ مطلوب؟}};

\end{tikzpicture}%
}
\end{center}

\section*{۸. جدول جامع تقابل‌ها}

\begin{table}[H]
\centering
\small
\begin{tabular}{|>{\columncolor{darkgray!10}}r|p{4cm}|p{4cm}|p{3.5cm}|}
\hline
\rowcolor{darkgray!30}
\textbf{تقابل} & \textbf{قطب چپ} & \textbf{قطب راست} & \textbf{میانه} \\
\hline
فرد/جمع & جمع‌گرایی & فردگرایی & توازن حقوق و وظایف \\
\hline
برابری/آزادی & برابری مقدم & آزادی مقدم & آزادی + فرصت برابر \\
\hline
دولت/بازار & دولت هدایت‌گر & بازار آزاد & اقتصاد مختلط \\
\hline
تغییر/سنت & تغییر بنیادین & حفظ سنت & اصلاح تدریجی \\
\hline
جهان/ملت & انترناسیونالیسم & ناسیونالیسم & همکاری + حاکمیت \\
\hline
آرمان/واقعیت & آرمان‌گرایی & واقع‌گرایی & واقع‌گرایی اخلاقی \\
\hline
خصوصی/عمومی & مالکیت اجتماعی & مالکیت خصوصی & اقتصاد مختلط \\
\hline
همکاری/رقابت & همکاری اجتماعی & رقابت بازار & رقابت تنظیم‌شده \\
\hline
بلندمدت/کوتاه‌مدت & برنامه‌ریزی بلندمدت & سود فوری & پایداری \\
\hline
\end{tabular}
\caption{خلاصهٔ تقابل‌های موضوعی}
\end{table}